%!TEX root = ../main.tex
% APÉNDICE B — Detalles técnicos e implementación
% (ocupa el segundo lugar en la lista de apéndices; ver main.tex)

\chapter{Detalles técnicos e implementación}\label{apx:tecnicos}

\begin{guia}
Todo lo que alguien necesita para reproducir los resultados:
hiperparámetros exactos, identificadores y versiones de los datos,
versiones de librerías, variables de entorno (sin secretos) y esquemas
de base de datos. Reproducibilidad total es un marcador de excelencia.
\end{guia}

\section{Configuración del modelo}

El modelo de gravedad en producción es un \texttt{RandomForestClassifier}
de scikit-learn, entrenado sobre 23.338 registros de siniestros con
motocicleta cruzados con clima histórico. La configuración de
hiperparámetros, ya fija, es la siguiente.

\begin{table}[H]
    \centering
    \caption{Hiperparámetros del modelo de gravedad en producción}
    \label{tab:hiperparametros}
    \small
    \begin{tabularx}{\textwidth}{l l L}
        \toprule
        \tb{Parámetro} & \tb{Valor} & \tb{Criterio} \\
        \midrule
        n\_estimators & 300 & cantidad de árboles combinados en el bosque \\
        max\_depth & 12 & profundidad máxima por árbol, evita que un árbol memorice casos particulares \\
        min\_samples\_leaf & 5 & mínimo de muestras por hoja, mismo objetivo que el límite anterior \\
        class\_weight & balanced\_subsample & compensa el desbalance entre las clases LEVE, GRAVE y MORTAL dándoles más peso relativo a las minoritarias en cada árbol \\
        random\_state & 42 & semilla fija, para que el entrenamiento sea reproducible \\
        \bottomrule
    \end{tabularx}
\end{table}

El modelo recibe 53 columnas de entrada después de codificación one-hot
(hora, temperatura, humedad, índice de inseguridad general por comuna,
día, comuna, tipo de vía, contraparte, sexo y condición de lluvia, entre
las variables de origen). La edad de la víctima y su indicador de edad
desconocida se excluyeron de la versión final: esa variante generalizó
peor entre años que la que no las incluye, según el protocolo de
evaluación temporal retrospectiva descrito en el
capítulo~\ref{cap:marco-tecnologico}. El artefacto entrenado no tiene
una tabla propia de versiones en la base de datos: se guarda
directamente como un archivo serializado (pickle) en el repositorio.

\section{Datos utilizados}

Los siniestros históricos de motociclistas que entrenan el modelo se
procesan fuera de línea, como archivos CSV, en un pipeline que no vive
en la base de datos de producción. No constituyen una tabla del esquema
relacional: su identificador, fuente y fecha de descarga quedan
documentados junto con el propio artefacto del modelo en el
repositorio, no en este apéndice, para no duplicar una versión que
pueda desactualizarse frente a la real.

\section{Arquitectura de la aplicación}

\subsection{Stack tecnológico}

\begin{itemize}
    \item \textbf{Interfaz}: React con Vite (aplicación de una sola página).
    \item \textbf{Mapa}: MapLibre GL JS.
    \item \textbf{API}: FastAPI (Python), con el modelo Random Forest cargado en memoria.
    \item \textbf{Base de datos}: PostgreSQL 17 con la extensión PostGIS, alojada en Supabase.
    \item \textbf{Modelo de gravedad}: Python, pandas y scikit-learn.
    \item \textbf{Contenedores}: Docker y Docker Compose, con servicios diferenciados para interfaz, API, base de datos y el proceso periódico de ingesta de prensa.
    \item \textbf{Despliegue}: Render (API y sitio web, como Web Services Docker del plan gratuito).
    \item \textbf{Tarea periódica}: scraper de prensa en GitHub Actions, con ejecución programada cada tres horas.
\end{itemize}

Las versiones exactas de cada dependencia (\texttt{requirements.txt} del
backend y del modelo, \texttt{package.json} del frontend) quedan fijadas
por archivo de bloqueo en el repositorio del proyecto, citado en el
anexo~\ref{anx:repositorios}; no se transcriben aquí para evitar que este
apéndice quede desactualizado frente al repositorio real.

\subsection{Esquema de base de datos}

El esquema real del backend en PostgreSQL/PostGIS está compuesto por las
tablas que se detallan en la Figura~\ref{fig:4-2} del
capítulo~\ref{cap:marco-tecnologico}: \texttt{hotspots\_riesgo} y
\texttt{corredores\_riesgo} (capas de riesgo vial derivadas por
agregación espacial), \texttt{hotspots\_inseguridad} (capa de contexto
de inseguridad general), \texttt{reportes\_prensa} (noticias de prensa
clasificadas), \texttt{reportes\_colaborativos} con
\texttt{reportes\_confirmaciones} (reportes cargados por usuarios y su
validación por otros usuarios), \texttt{usuarios} y
\texttt{codigos\_auth} (esquema de cuentas y verificación por código).
Las tablas \texttt{hotspots\_riesgo} y \texttt{hotspots\_inseguridad}
guardan sus coordenadas también como geometría nativa de PostGIS, aunque
el backend hoy calcula proximidad con la fórmula de Haversine en código
Python en lugar de una consulta espacial.

\subsection{Variables de entorno}

\begin{table}[H]
    \centering
    \caption{Variables de entorno declaradas por la aplicación (sin valores reales)}
    \label{tab:env}
    \small
    \begin{tabularx}{\textwidth}{l L}
        \toprule
        \tb{Variable} & \tb{Propósito} \\
        \midrule
        DATABASE\_URL & cadena de conexión a PostgreSQL/PostGIS en Supabase (pooler de sesión) \\
        GMAIL\_OAUTH\_TOKEN & token de larga duración para el envío de códigos de verificación por la API HTTP de Gmail \\
        SESSION\_SECRET & clave usada para firmar la cookie de sesión (HMAC) \\
        MODEL\_PATH & ruta del artefacto serializado (pickle) del modelo de gravedad \\
        \bottomrule
    \end{tabularx}
\end{table}

Ningún secreto real (contraseñas, tokens, claves) se versiona en el
repositorio: las variables de entorno se cargan en el panel de cada
proveedor (Render, Supabase, GitHub Actions) y el repositorio incluye
únicamente un archivo de ejemplo con los nombres de las variables, sin
valores.
